\section{Syntax at the Word Level}
\label{sec:word-level-syntax}
\slidesrc{Syntax: RE \& word: 16--82}

Before a sentence can be tagged or parsed, the text must be cut into words, and the
words must often be cleaned and normalised. This section covers the word-level tasks of
the syntactic level (Section~\ref{sec:course-map}):
\begin{enumerate}[nosep]
  \item \textbf{word segmentation}: where does a word start and end?
        (Section~\ref{sec:word-segmentation});
  \item \textbf{stop words} and \textbf{TF-IDF}: which words carry little information,
        and how to measure how important a word is for a document
        (Section~\ref{sec:stop-words-tfidf});
  \item \textbf{stemming and lemmatisation}: reducing \emph{organizes},
        \emph{organizing} to a common form, with the Porter stemmer in detail
        (Section~\ref{sec:stemming-lemmatization});
  \item \textbf{minimum edit distance}: how similar are two strings, computed with
        dynamic programming (Section~\ref{sec:minimum-edit-distance}).
\end{enumerate}

\subsection{Word Segmentation}
\label{sec:word-segmentation}
\slidesrc{Syntax: RE \& word: 16--20}

\begin{nbdef}{Word segmentation}
\textbf{Word segmentation} (also \emph{tokenisation}) is the problem of dividing a string
of written language into its component words.
\end{nbdef}

\subsubsection{Why it is not trivial}
\slidesrc{Syntax: RE \& word: 17}

In English and in many other languages written with the Latin alphabet, the
\textbf{space} is a good approximation of a word divider (word delimiter). It is only an
approximation, because languages vary in how they treat collocations and compounds
(Section~\ref{sec:morphology}): many English compound nouns are written in several ways,
e.g.\ \emph{multi-agent} = \emph{multiagent} = \emph{multi agent}, and speakers differ on
whether they think of them as one noun or as a noun phrase. Punctuation adds more cases
(\emph{don't}, \emph{U.S.A.}, \emph{3.14}, \emph{e-mail}).

Many scripts have no word space at all, and there word segmentation is a hard problem:
\begin{itemize}[nosep]
  \item \textbf{Chinese} and \textbf{Japanese}: sentences are delimited, words are not;
  \item \textbf{Thai} and \textbf{Lao}: phrases and sentences are delimited, words are not;
  \item \textbf{Vietnamese}: syllables are delimited, words are not.
\end{itemize}

\subsubsection{The MaxMatch algorithm}
\slidesrc{Syntax: RE \& word: 18}

A simple baseline for Chinese is the greedy \textbf{maximum matching} algorithm. It needs
a \textbf{dictionary} (word list) of the language.

\begin{nbdef}{MaxMatch}
\begin{enumerate}[nosep]
  \item Put a pointer at the beginning of the string.
  \item Choose the \textbf{longest} dictionary word that matches the input starting at
        the pointer, and move the pointer to the end of that word.
  \item If no word matches, move the pointer one character forward instead (creating a
        one-character ``word'').
  \item Repeat from the new pointer position until the end of the string.
\end{enumerate}
\end{nbdef}

\begin{nbexample}{MaxMatch on English without spaces}
Input: \texttt{wecanonlyseeashortdistanceahead} (``we can only see a short distance
ahead'').
\begin{center}
\small
\begin{tabular}{@{}l l l@{}}
\toprule
\textbf{Remaining input} & \textbf{Longest word} & \textbf{Why} \\
\midrule
\texttt{wecanonly\ldots}  & \emph{we}       & \\
\texttt{canonlysee\ldots} & \emph{canon}    & longer than \emph{can} \\
\texttt{lyseea\ldots}     & \emph{l}        & no word starts here: one character \\
\texttt{yseea\ldots}      & \emph{y}        & same \\
\texttt{seeashort\ldots}  & \emph{see}      & \\
\texttt{ashort\ldots}     & \emph{ash}      & longer than \emph{a} and \emph{as} \\
\texttt{ort\ldots}        & \emph{o}, \emph{r}, \emph{t} & three single characters \\
\texttt{distanceahead}    & \emph{distance}, \emph{ahead} & \\
\bottomrule
\end{tabular}
\end{center}
Output: \emph{we canon l y see ash o r t distance ahead}.
\end{nbexample}

\begin{nbintuition}[why greedy fails here and works for Chinese]
Greedy choices are never undone: once \emph{canon} is taken, the correct \emph{can only}
is impossible. In English, words are long and many short words hide inside longer ones,
so this happens constantly. Chinese words are mostly one or two characters long, so the
longest match is usually right, and MaxMatch is a surprisingly strong baseline there.
Better segmenters score \emph{whole} segmentations (e.g.\ with word probabilities)
instead of committing greedily.
\end{nbintuition}

\subsubsection{NLTK}
\slidesrc{Syntax: RE \& word: 19--20}

The course uses \textbf{NLTK} (Natural Language Toolkit, \url{https://www.nltk.org/})
to show some problems and some of their possible solutions. NLTK is a Python platform
for working with human language data: it provides interfaces to over 50 corpora and
lexical resources such as WordNet (Section~\ref{sec:wordnet}), and libraries for
classification, tokenisation, stemming, tagging, parsing and semantic reasoning. The
course also uses less known, but elegant, tools.

\begin{minted}{python}
import nltk
# nltk.download()   # only once, to download corpora and models
from nltk.tokenize import word_tokenize
print(word_tokenize("Don't segment U.S.A. badly, Mr. Smith!"))
# ['Do', "n't", 'segment', 'U.S.A.', 'badly', ',', 'Mr.', 'Smith', '!']
\end{minted}

Note how the tokeniser splits \emph{don't} into \emph{Do} + \emph{n't} (two
morphemes, Section~\ref{sec:morphology}) but keeps \emph{U.S.A.} and \emph{Mr.}
together: the full stop is not always a word boundary.

\subsection{Stop Words and TF-IDF}
\label{sec:stop-words-tfidf}
\slidesrc{Syntax: RE \& word: 21--47}

\subsubsection{Stop words}
\slidesrc{Syntax: RE \& word: 21--23}

\begin{nbdef}{Stop words}
\textbf{Stop words} are words that are filtered out before or after processing natural
language text. The term usually refers to the \emph{most common} words of a language
(\emph{the}, \emph{of}, \emph{and}, \emph{a}, \emph{to}, \ldots), but there is no single
universal list used by all NLP tools, and some tools do not use one at all.
\end{nbdef}

\begin{itemize}[nosep]
  \item \textbf{How a stop list is built}: sort the terms by \emph{collection frequency}
        (the total number of times each term appears in the document collection), take
        the most frequent ones, and often filter them by hand according to their meaning
        in the domain of the documents. The words of the stop list are then discarded
        during indexing. Example lists: \url{https://www.ranks.nl/stopwords}.
  \item \textbf{When not to remove them}: some tools avoid removing stop words to
        support \emph{phrase search} (exact, not keyword-based): ``to be or not to be''
        consists almost only of stop words.
\end{itemize}

A word cloud of a novel (the slides show one for a Thomas Hardy novel) is dominated by
\emph{the}, \emph{and}, \emph{of}, \emph{to}, \emph{a}, \emph{was}, \emph{her}: they are
frequent in \emph{any} English text and say nothing about this one.

\subsubsection{Frequent in the document or frequent in the language?}
\slidesrc{Syntax: RE \& word: 24--29}

\begin{nbintuition}[the key question]
Is a word that appears often in a document always a stop word? No. A word can be
frequent in a document $d$ for two very different reasons:
\begin{itemize}[nosep]
  \item it is related to the \textbf{topic} of $d$ (``cat'' in a document about cats);
  \item it is frequent in the \textbf{language} (``the'' in every document).
\end{itemize}
To tell them apart, we must look not only at $d$ but at a whole \textbf{corpus}: a
topical word is frequent in \emph{some} documents, a stop word in \emph{all} of them.
\end{nbintuition}

The running example uses two short documents:

\begin{nbexample}{the corpus}
\textbf{Doc1} (54 words): \emph{The cat (Felis catus) is a small carnivorous mammal. It
is the only domesticated species in the family Felidae and often referred to as the
domestic cat to distinguish it from wild members of the family. The cat is either a
house cat or a farm cat, which are pets, or a feral cat.}

\textbf{Doc2} (42 words): \emph{A car (or automobile) is a wheeled motor vehicle used for
transportation. Cars became widely available in the early 20th century. One of the first
cars accessible to the masses was the 1908 Model T. Cars were rapidly adopted in the US.}

\begin{center}
\small
\begin{tabular}{@{}l c c c@{}}
\toprule
 & \textbf{``the''} & \textbf{``cat''} & \textbf{``car(s)''} \\
\midrule
Doc1 & 6 & 6 & 0 \\
Doc2 & 5 & 0 & 4 \\
\midrule
corpus $C_2 = \{\text{Doc1}, \text{Doc2}\}$ & 11 & 6 & 4 \\
\bottomrule
\end{tabular}
\end{center}
Corpus $C_1$ contains only Doc1; corpus $C_2$ contains Doc1 and Doc2.
\end{nbexample}

\begin{nbwarning}[a counting slip in the slides]
The slides count \textbf{4} occurrences of ``the'' in Doc2 (and 10 in $C_2$). Doc2
actually contains it \textbf{5} times: \emph{\underline{the} early}, \emph{\underline{the}
first}, \emph{\underline{the} masses}, \emph{\underline{the} 1908}, \emph{\underline{the}
US}. This changes $\text{tf}(\text{``the''}, \text{Doc2})$ from $4/42$ to $5/42$, but not
the final TF-IDF value, which is $0$ either way (see below).
\end{nbwarning}

In $C_1$ alone, ``the'' and ``cat'' are equally frequent (6 times each): with a single
document there is no way to tell a topical word from a stop word. In $C_2$ the
difference appears: ``the'' is in \emph{both} documents, ``cat'' only in one.

\subsubsection{Term frequency and inverse document frequency}
\slidesrc{Syntax: RE \& word: 30--34}

\begin{nbdef}{TF-IDF}
\textbf{Term frequency}: the number of times term $t$ occurs in document $d$, usually
normalised by the length of the document:
\[
  \text{tf}(t, d) = \frac{\text{number of occurrences of } t \text{ in } d}{\text{length}(d)}.
\]
\textbf{Inverse document frequency}: the logarithm of the inverse fraction of documents
of the corpus $C$ that contain $t$:
\[
  \text{idf}(t, C) = \log \frac{\abs{C}}{\abs{\{d \in C : t \in d\}}},
\]
where $\abs{C}$ is the number of documents in the corpus.

\textbf{TF-IDF}: the product of the two,
\[
  \text{tf-idf}(t, d, C) = \text{tf}(t, d) \cdot \text{idf}(t, C).
\]
\end{nbdef}

\begin{itemize}
  \item \textbf{tf} is high when $t$ is frequent \emph{in this document}.
  \item \textbf{idf} is high when $t$ is \emph{rare in the corpus}. If $t$ occurs in
        every document, the fraction is $\abs{C}/\abs{C} = 1$ and $\text{idf} = \log 1 = 0$:
        the word cannot distinguish documents at all, exactly like a stop word.
  \item The \textbf{logarithm} dampens the effect: a word in 1 document out of 1000 is
        more informative than one in 10 out of 1000, but not 10 times more.
  \item If $t$ appears in \emph{no} document the denominator is 0. A common fix is to
        add 1 to the denominator; it makes sense when the corpus is large enough, since
        it slightly distorts the values.
\end{itemize}

\begin{nbwarning}[base of the logarithm]
The slides' numbers use the base-10 logarithm ($\log 2 \approx 0.301$). Other sources use
the natural logarithm ($\ln 2 \approx 0.693$). The base only rescales all idf values by
the same factor, so the \emph{ranking} of words does not change; just be consistent.
\end{nbwarning}

\subsubsection{Worked example on \texorpdfstring{$C_2$}{C2}}
\slidesrc{Syntax: RE \& word: 35--46}

\begin{nbexample}{TF-IDF of ``the'', ``cat'', ``car''}
First the idf values ($\abs{C_2} = 2$):
\begin{align*}
  \text{idf}(\text{the}, C_2) &= \log\tfrac{2}{2} = \log 1 = 0
      \reason{``the'' is in both documents}\\
  \text{idf}(\text{cat}, C_2) &= \log\tfrac{2}{1} = \log 2 \approx 0.301
      \reason{only in Doc1}\\
  \text{idf}(\text{car}, C_2) &= \log\tfrac{2}{1} = \log 2 \approx 0.301
      \reason{only in Doc2}
\end{align*}
Then the products:
\begin{align*}
  \text{tf-idf}(\text{the}, \text{Doc1}, C_2) &= \tfrac{6}{54}\cdot 0 = 0 \\
  \text{tf-idf}(\text{the}, \text{Doc2}, C_2) &= \tfrac{5}{42}\cdot 0 = 0 \\
  \text{tf-idf}(\text{cat}, \text{Doc1}, C_2) &= \tfrac{6}{54}\cdot 0.301 \approx 0.111 \cdot 0.301 \approx 0.033 \\
  \text{tf-idf}(\text{cat}, \text{Doc2}, C_2) &= 0 \cdot 0.301 = 0
      \reason{tf $=0$, idf irrelevant}\\
  \text{tf-idf}(\text{car}, \text{Doc1}, C_2) &= 0 \cdot 0.301 = 0 \\
  \text{tf-idf}(\text{car}, \text{Doc2}, C_2) &= \tfrac{4}{42}\cdot 0.301 \approx 0.095 \cdot 0.301 \approx 0.029
\end{align*}
\end{nbexample}

\begin{center}
\small
\begin{tabular}{@{}l c c c@{}}
\toprule
\textbf{tf-idf} & \textbf{``the''} & \textbf{``cat''} & \textbf{``car''} \\
\midrule
Doc1 & 0 & \textbf{0.033} & 0 \\
Doc2 & 0 & 0 & \textbf{0.029} \\
\bottomrule
\end{tabular}
\captionof{table}{TF-IDF in $C_2$. Each document's only non-zero word is its topic.}
\end{center}

``The'' is as frequent as ``cat'' in Doc1, but its TF-IDF is $0$ because it appears
everywhere; ``cat'' and ``car'' come out as the words that characterise their documents.
This is how TF-IDF finds \emph{keywords} and, symmetrically, \emph{stop words} (terms with
idf close to 0) automatically. Online analysers such as
\url{https://www.online-utility.org/text/analyzer.jsp} compute word frequencies of a
text you paste.

\begin{nbwhy}
TF-IDF is one of the text features of Section~\ref{sec:course-map}: representing each
document by the TF-IDF values of its words gives a vector that search engines use to
rank documents for a query and classifiers use to decide the topic of a text.
\end{nbwhy}

\subsection{Stemming and Lemmatisation}
\label{sec:stemming-lemmatization}
\slidesrc{Syntax: RE \& word: 48--65}

\subsubsection{The goal: one form for many}
\slidesrc{Syntax: RE \& word: 48}

Documents use different forms of a word: \emph{organize}, \emph{organizes},
\emph{organizing} (inflectional morphemes, Section~\ref{sec:morphology}); and there are
families of derivationally related words with similar meanings: \emph{democratic},
\emph{democratization} (derivational morphemes). For a search, it is often useful that
a query for one of these words also returns documents containing another word of the
same set.

The goal of both stemming and lemmatisation is to \textbf{reduce inflectional forms, and
sometimes derivationally related forms, of a word to a common base form}:
\begin{align*}
  \text{am, are, is} &\;\to\; \text{be}\\
  \text{car, cars, car's, cars'} &\;\to\; \text{car}\\
  \text{the boy's cars are different colors} &\;\to\; \text{the boy car be differ color}
\end{align*}

\subsubsection{Stemming vs.\ lemmatisation}
\slidesrc{Syntax: RE \& word: 49}

\begin{nbdef}{Stemming and lemmatisation}
\begin{itemize}[nosep]
  \item \textbf{Stemming} is a \emph{crude heuristic} process that chops off the ends of
        words, hoping to get it right most of the time. It often removes derivational
        affixes too. The result, the \emph{stem}, need not be a real word.
  \item \textbf{Lemmatisation} does things properly, with a \emph{vocabulary} and a
        \emph{morphological analysis} of words. It normally removes inflectional endings
        only and returns the base or dictionary form of the word, the \textbf{lemma}.
\end{itemize}
Stemming most commonly collapses derivationally related words; lemmatisation commonly
collapses only the different inflectional forms of a lemma.
\end{nbdef}

\begin{center}
\small
\begin{tabular}{@{}l l l@{}}
\toprule
\textbf{Word} & \textbf{Porter stem} & \textbf{Lemma} \\
\midrule
are          & ar       & be \\
theses       & these    & thesis \\
universal    & univers  & universal \\
university   & univers  & university \\
generalizations & gener & generalization \\
\bottomrule
\end{tabular}
\captionof{table}{A stemmer chops, a lemmatiser looks words up. Stems are not words, and
unrelated words can end up with the same stem.}
\end{center}

\subsubsection{The Porter stemmer}
\slidesrc{Syntax: RE \& word: 50--52}

The \textbf{Porter stemmer} (Martin Porter, 1980) is the classic stemming algorithm for
English. It is a simple \textbf{rule-based} and \textbf{heuristic} method: a cascade of
suffix-rewriting rules such as \textsc{ational} $\to$ \textsc{ate}
(\emph{relational} $\to$ \emph{relate}), applied in sets, in order. The slides present
it in seven steps.

\paragraph{Consonants, vowels and the measure.}
\begin{itemize}[nosep]
  \item A \textbf{consonant} is a letter other than \texttt{a}, \texttt{e}, \texttt{i},
        \texttt{o}, \texttt{u}, and other than \texttt{y} when preceded by a consonant.
        A \textbf{vowel} is any other letter. (So the \texttt{y} of \emph{toy} is a
        consonant, the \texttt{y} of \emph{happy} is a vowel.)
  \item Write $C$ for a sequence of one or more consonants and $V$ for a sequence of one
        or more vowels. Then \emph{every} word has the form
        \[ [C]\,(VC)^{m}\,[V] \]
        where the square brackets mean ``optional''. The number $m$ is the
        \textbf{measure} of the word (or stem): roughly, its number of syllables.
\end{itemize}

\begin{nbexample}{measures}
\begin{center}
\small
\begin{tabular}{@{}l l l c@{}}
\toprule
\textbf{Word} & \textbf{Letters as C/V} & \textbf{Form} & $m$ \\
\midrule
tree     & tr$\cdot$ee                        & $C\,V$                & 0 \\
trouble  & tr$\cdot$ou$\cdot$bl$\cdot$e       & $C\,(VC)\,V$          & 1 \\
troubles & tr$\cdot$ou$\cdot$bl$\cdot$e$\cdot$s & $C\,(VC)(VC)$       & 2 \\
oaten    & oa$\cdot$t$\cdot$e$\cdot$n         & $(VC)(VC)$            & 2 \\
\bottomrule
\end{tabular}
\end{center}
The slides write \emph{troubles} as ``C V CVC'': grouping it as $C\,(VC)(VC)$ gives
$m = 2$.
\end{nbexample}

\paragraph{Rule format.} Each rule has the form \quad(\emph{condition})\; $S_1 \to S_2$
\quad: if the word ends with suffix $S_1$ and the \emph{stem} (what is left after
removing $S_1$) satisfies the condition, $S_1$ is replaced by $S_2$. The conditions are:

\begin{center}
\small
\begin{tabular}{@{}l l@{}}
\toprule
\textbf{Condition} & \textbf{Meaning} \\
\midrule
$m > k$, $m = k$ & the measure of the stem \\
\texttt{*S}  & the stem ends with the letter \texttt{S} (similarly \texttt{*L}, \texttt{*Z}) \\
\texttt{*v*} & the stem contains a vowel \\
\texttt{*d}  & the stem ends with a double consonant (e.g.\ \texttt{-tt}, \texttt{-ss}) \\
\texttt{*o}  & the stem ends with $cvc$, where the second $c$ is not \texttt{w}, \texttt{x}, \texttt{y} \\
\bottomrule
\end{tabular}
\end{center}

The measure condition is what makes the stemmer safe on short words: (\emph{m}$>$0)
\textsc{ness} $\to \varepsilon$ turns \emph{goodness} into \emph{good}, but it would not
touch a hypothetical word whose stem would be empty.

\subsubsection{The seven steps}
\slidesrc{Syntax: RE \& word: 53--58}

In the rules, $\varepsilon$ is the empty string (the suffix is deleted).

\begin{center}
\footnotesize
\begin{tabularx}{\linewidth}{@{}>{\raggedright\arraybackslash}p{1.9cm} >{\raggedright\arraybackslash}X >{\raggedright\arraybackslash}X@{}}
\toprule
\textbf{Step} & \textbf{Rules} & \textbf{Examples} \\
\midrule
1 (plurals) & \textsc{sses}$\to$\textsc{ss};\; \textsc{ies}$\to$\textsc{i};\;
  \textsc{ss}$\to$\textsc{ss};\; \textsc{s}$\to\varepsilon$ &
  caresses$\to$caress, ponies$\to$poni, ties$\to$ti, caress$\to$caress, cats$\to$cat \\
\addlinespace
2a (\emph{-ed}, \emph{-ing}) & ($m{>}0$) \textsc{eed}$\to$\textsc{ee};\;
  (\texttt{*v*}) \textsc{ed}$\to\varepsilon$;\; (\texttt{*v*}) \textsc{ing}$\to\varepsilon$ &
  agreed$\to$agree, feed$\to$feed; plastered$\to$plaster, bled$\to$bled;
  motoring$\to$motor, sing$\to$sing \\
\addlinespace
2b (clean-up, only if the 2nd or 3rd rule of 2a fired) &
  \textsc{at}$\to$\textsc{ate};\; \textsc{bl}$\to$\textsc{ble};\;
  (\texttt{*d} and not \texttt{*L}, \texttt{*S}, \texttt{*Z}) $\to$ single letter;\;
  ($m{=}1$ and \texttt{*o}) $\to$ \textsc{e} &
  conflat(ed)$\to$conflate, troubl(ing)$\to$trouble, hopp(ing)$\to$hop,
  tann(ed)$\to$tan, fall(ing)$\to$fall, fil(ing)$\to$file, fail$\to$fail \\
\addlinespace
3 (\emph{y}) & (\texttt{*v*}) \textsc{y}$\to$\textsc{i} & happy$\to$happi, sky$\to$sky \\
\addlinespace
4 (derivational I) & ($m{>}0$) \textsc{ational}$\to$\textsc{ate};\;
  \textsc{ization}$\to$\textsc{ize};\; \textsc{biliti}$\to$\textsc{ble};\; \ldots &
  relational$\to$relate, generalization$\to$generalize, sensibiliti$\to$sensible \\
\addlinespace
5 (derivational II) & ($m{>}0$) \textsc{icate}$\to$\textsc{ic};\;
  \textsc{ful}$\to\varepsilon$;\; \textsc{ness}$\to\varepsilon$;\; \ldots &
  triplicate$\to$triplic, hopeful$\to$hope, goodness$\to$good \\
\addlinespace
6 (derivational III) & ($m{>}1$) \textsc{ance}, \textsc{ent}, \textsc{ive},
  \textsc{er}, \textsc{al}, \textsc{ant}, \ldots $\to\varepsilon$ &
  allowance$\to$allow, dependent$\to$depend, effective$\to$effect \\
\addlinespace
7a (final \emph{e}) & ($m{>}1$) \textsc{e}$\to\varepsilon$;\;
  ($m{=}1$ and not \texttt{*o}) \textsc{e}$\to\varepsilon$ & probate$\to$probat,
  cease$\to$ceas \\
\addlinespace
7b (double \emph{l}) & ($m{>}1$ and \texttt{*d} and \texttt{*L}) $\to$ single letter &
  controll$\to$control, roll$\to$roll \\
\bottomrule
\end{tabularx}
\captionof{table}{The Porter stemmer as presented in the slides (in Porter's paper the
same rules are numbered 1a, 1b, 1c, 2, 3, 4, 5a, 5b).}
\end{center}

\begin{nbwarning}[corrections to the Porter slides]
\begin{itemize}[nosep]
  \item \textbf{Consonant definition}: the slide reads ``a letter other than A, E, I, O,
        U, and Y preceded by consonant''. Porter's definition is ``other than A, E, I, O,
        U, \emph{and other than} Y preceded by a consonant'': a Y after a consonant is a
        vowel (\emph{happy}), otherwise a consonant (\emph{toy}).
  \item \textbf{Step 2a}: the slide writes the condition of \textsc{eed} as $m>1$, then
        uses \emph{agreed} $\to$ \emph{agree} as an example, but \emph{agr} has $m=1$.
        Porter's rule is $(m>0)$ \textsc{eed} $\to$ \textsc{ee}, as in the table above.
  \item \textbf{Step 7a}: the slide's second rule is ``($m{=}1$ and not \texttt{*o})
        \textsc{ness} $\to \varepsilon$'' with \emph{goodness} as the example (copied
        from step 5). In Porter's algorithm this rule removes a final \textsc{e}:
        \emph{cease} $\to$ \emph{ceas}.
\end{itemize}
\end{nbwarning}

\subsubsection{Examples and problems}
\slidesrc{Syntax: RE \& word: 59--62}

\begin{nbexample}{running the steps}
\begin{itemize}[nosep]
  \item \emph{controlling} $\xrightarrow{\text{2a: } -\textsc{ing}}$ \emph{controll}
        $\xrightarrow{\text{7b}}$ \emph{control}.
  \item \emph{singing} $\xrightarrow{\text{2a: } -\textsc{ing}}$ \emph{sing} (no further
        rule applies).
  \item \emph{computers} $\xrightarrow{\text{1: } -\textsc{s}}$ \emph{computer}
        $\xrightarrow{\text{6: } -\textsc{er}}$ \emph{comput}.
  \item \emph{generalizations} $\xrightarrow{\text{1}}$ \emph{generalization}
        $\xrightarrow{\text{4: \textsc{ization}}\to\textsc{ize}}$ \emph{generalize}
        $\xrightarrow{\text{5: \textsc{alize}}\to\textsc{al}}$ \emph{general}
        $\xrightarrow{\text{6: } -\textsc{al}}$ \emph{gener}.
\end{itemize}
\end{nbexample}

\begin{nbwarning}[the examples slide]
On the slide, \emph{singing} ends with ``Step 6, Rule 4 $\to$ compute'', a line copied
from the \emph{computers} example: \emph{singing} simply becomes \emph{sing}. The slide
also stops \emph{computers} at \emph{compute} and \emph{generalizations} at
\emph{general}. Running the complete algorithm (e.g.\ NLTK's implementation) gives
\emph{comput} and \emph{gener}: stems are not required to be words, and Porter's
algorithm removes \textsc{er} and \textsc{al} in its last derivational step.
\end{nbwarning}

\paragraph{Problems.} Being heuristic, the stemmer makes two kinds of errors:
\begin{itemize}[nosep]
  \item \textbf{over-stemming}: unrelated words get the same stem. \emph{elephants}
        $\to$ \emph{elephant} (step 1) $\to$ \emph{eleph} (step 6 removes \textsc{ant} as
        if it were a suffix); \emph{university} and \emph{universal} both become
        \emph{univers};
  \item \textbf{under-stemming}: forms of the same word get different stems.
        \emph{thesis} $\to$ \emph{thesi} but \emph{theses} $\to$ \emph{these}.
\end{itemize}
The slides also list \emph{doing} $\to$ \emph{doe} as a problem; the standard algorithm
actually returns \emph{do} for \emph{doing} (the \textsc{e} of step 2b is added only when
$m=1$, and \emph{do} has $m=0$).

References: Porter, M.\,F.\ (1980), \emph{An algorithm for suffix stripping},
Program 14(3):130--137; the Porter stemmer home page
\url{https://tartarus.org/martin/PorterStemmer/}; Jurafsky and Martin, Section~3.4.
Online demos: \url{https://9ol.es/porter_js_demo.html}; stemmers for other languages:
\url{http://text-processing.com/demo/stem/}.

\subsubsection{Two texts to test a stemmer}
\slidesrc{Syntax: RE \& word: 63--65}

The slides propose two short texts full of traps. Here is what a real stemmer does with
them.

\begin{nbexample}{Italian: a Snowball stemmer}
\emph{Sono stato in uno stato il cui statuto prevede che la biscia non possa entrare in
una bisca. Una messe di polemiche è piovuta sullo statuto, per le limitazioni messe
alla biscia, che è andata a messa insieme al console per consolarsi.}
\begin{center}
\small
\begin{tabular}{@{}l l l@{}}
\toprule
\textbf{Words} & \textbf{Meanings} & \textbf{Stem} \\
\midrule
stato, stato, statuto & ``been'', ``state'', ``statute'' & \emph{stat} \\
biscia, bisca & ``grass snake'', ``gambling den'' & \emph{bisc} \\
messe, messe, messa & ``harvest'', ``put'' (participle), ``mass'' & \emph{mess} \\
console, consolarsi & ``consul'', ``to console oneself'' & \emph{consol} \\
\bottomrule
\end{tabular}
\end{center}
Every group is over-stemmed: completely unrelated words collapse. Even a lemmatiser
needs context here: the two \emph{stato} have different lemmas (\emph{essere}, the verb;
\emph{stato}, the noun), and only the part of speech tells them apart
(Section~\ref{sec:pos-tagging}).
\end{nbexample}

\begin{nbexample}{English: Miss Marple}
\emph{Once Miss Marple had a cold and she went to the doctor: ``I go!'', she said, and she
read the doctor's doctoral thesis. ``These are good theses!'', and -- although she missed
her train -- Miss Marple understood that University has a universal mission.}

Porter gives: \emph{Miss}, \emph{missed} $\to$ \emph{miss} (a title and a verb merged);
\emph{doctor}, \emph{doctoral} $\to$ \emph{doctor}; \emph{University}, \emph{universal}
$\to$ \emph{univers}; \emph{theses} $\to$ \emph{these}, the same as the pronoun
\emph{These}, while \emph{thesis} $\to$ \emph{thesi}; \emph{has} $\to$ \emph{ha};
\emph{are} $\to$ \emph{ar}. A lemmatiser instead returns \emph{thesis} for
\emph{theses}, \emph{be} for \emph{are}, \emph{have} for \emph{has}, and keeps
\emph{university} and \emph{universal} apart.
\end{nbexample}

\begin{nbintuition}[which one to use]
Stemming is fast, needs no dictionary and increases \emph{recall} in search (more forms
of a word are found), at the price of \emph{precision} (unrelated words are merged).
Lemmatisation is slower and needs a vocabulary and often the part of speech, but its
output is made of real words. Choose according to which kind of error hurts your
application more.
\end{nbintuition}

\subsection{Minimum Edit Distance}
\label{sec:minimum-edit-distance}
\slidesrc{Syntax: RE \& word: 66--82}

\subsubsection{How similar are two strings?}
\slidesrc{Syntax: RE \& word: 66--68, 72--73}

Many tasks need a measure of how close two strings are:
\begin{itemize}
  \item \textbf{Spell correction}: the user typed \emph{graffe}. Which is closest:
        \emph{graf}, \emph{graft}, \emph{grail} or \emph{giraffe}?
  \item \textbf{Computational biology}: align two sequences of nucleotides (the bases
        Adenine, Guanine, Thymine/Uracil and Cytosine of DNA/RNA), giving each letter a
        partner letter or a gap \texttt{-}:
\begin{center}
\small\ttfamily
\begin{tabular}{@{}l@{}}
AGGCTATCACCTGACCTCCAGGCCGATGCCC\\
TAGCTATCACGACCGCGGTCGATTTGCCCGAC\\[3pt]
\textrm{\normalsize aligned:}\\
-AGGCTATCACCTGACCTCCAGGCCGA--TGCCC---\\
TAG-CTATCAC--GACCGC--GGTCGATTTGCCCGAC
\end{tabular}
\end{center}
        Alignments are used in molecular phylogenetics (evolutionary relationships from
        hereditary differences in DNA), in comparing the genomes of ill and healthy
        individuals, and to understand the function of unknown proteins.
  \item \textbf{Evaluating machine translation and speech recognition}: compare the
        system's output (hypothesis H) with a reference (R), counting substitutions (S),
        insertions (I) and deletions (D):
\begin{center}
\small
\begin{tabular}{@{}l *{9}{c}@{}}
R: & Spokesman & confirms & *   & senior & government & adviser & was & shot & *    \\
H: & Spokesman & said     & the & senior & *          & adviser & was & shot & dead \\
   &           & S        & I   &        & D          &         &     &      & I
\end{tabular}
\end{center}
  \item \textbf{Named entity extraction and coreference}: deciding that \emph{IBM Inc.}
        and \emph{IBM}, or \emph{Stanford President John Hennessy} and \emph{Stanford
        University President John Hennessy}, refer to the same entity.
\end{itemize}

\subsubsection{Definition}
\slidesrc{Syntax: RE \& word: 69--71}

\begin{nbdef}{Minimum edit distance}
The \textbf{minimum edit distance} between two strings is the minimum number of
editing operations
\begin{itemize}[nosep]
  \item \textbf{insertion} of a character,
  \item \textbf{deletion} of a character,
  \item \textbf{substitution} of a character with another,
\end{itemize}
needed to transform one string into the other. If operations have \emph{costs}, it is the
minimum total cost. With cost 1 for insertions and deletions and \textbf{cost 2 for
substitutions} it is called the \textbf{Levenshtein distance} (the version used in these
notes unless stated otherwise).
\end{nbdef}

An \textbf{alignment} shows the operations explicitly. For \emph{intention} and
\emph{execution} (\texttt{*} = gap):
\[
\begin{array}{*{10}{c}}
  \texttt{I} & \texttt{N} & \texttt{T} & \texttt{E} & \texttt{*} & \texttt{N} & \texttt{T} & \texttt{I} & \texttt{O} & \texttt{N}\\
  | & | & | & | & | & | & | & | & | & |\\
  \texttt{*} & \texttt{E} & \texttt{X} & \texttt{E} & \texttt{C} & \texttt{U} & \texttt{T} & \texttt{I} & \texttt{O} & \texttt{N}\\
  \text{d} & \text{s} & \text{s} & & \text{i} & \text{s} & & & &
\end{array}
\]
One deletion (d), three substitutions (s) and one insertion (i):
\begin{itemize}[nosep]
  \item if every operation costs 1, the distance is $1 + 3 + 1 = 5$;
  \item with Levenshtein costs (substitution $= 2$), it is $1 + 3\cdot 2 + 1 = 8$.
\end{itemize}

\begin{nbexample}{the spelling question}
Levenshtein distances from \emph{graffe}: \emph{giraffe} $=1$ (insert \texttt{i}),
\emph{graf} $=2$ (delete \texttt{f} and \texttt{e}), \emph{graft} $=3$ (delete
\texttt{f}, substitute \texttt{e}$\to$\texttt{t}), \emph{grail} $=5$. So \emph{giraffe}
is the best correction. (With all costs equal to 1: $1, 2, 2, 3$.)
\end{nbexample}

\subsubsection{Edit distance as search}
\slidesrc{Syntax: RE \& word: 74--78}

Finding the minimum edit distance is a search for a path (a sequence of edits) from the
start string to the final string:
\begin{itemize}[nosep]
  \item \textbf{initial state}: the word we are transforming;
  \item \textbf{operators}: insert, delete, substitute;
  \item \textbf{goal state}: the word we are trying to reach;
  \item \textbf{path cost}: what we minimise, the number (or cost) of edits.
\end{itemize}
The space of all edit sequences is huge, so we cannot explore it naively. But many
distinct paths end up in the \emph{same} state (the same prefix of $X$ turned into the
same prefix of $Y$), and for each such state we only need to remember the
\emph{cheapest} way to reach it. That is exactly the idea of dynamic programming.

\begin{nbdef}{The subproblems}
For strings $X$ of length $n$ and $Y$ of length $m$, let
\[
  D(i, j) = \text{edit distance between } X[1..i] \text{ and } Y[1..j],
\]
i.e.\ between the first $i$ characters of $X$ and the first $j$ characters of $Y$. The
edit distance between $X$ and $Y$ is $D(n, m)$.
\end{nbdef}

\textbf{Dynamic programming} computes $D(n,m)$ \emph{bottom-up} in a table, solving the
problem by combining solutions to subproblems: first $D(i,j)$ for small $i, j$, then
larger ones from the smaller values already computed.

\subsubsection{The algorithm}
\slidesrc{Syntax: RE \& word: 79--80}

\begin{nbfact}{Levenshtein distance by dynamic programming}
\textbf{Initialisation} (distance from the empty string):
\[ D(i, 0) = i, \qquad D(0, j) = j. \]
\textbf{Recurrence}: for each $i = 1, \dots, n$ and each $j = 1, \dots, m$,
\[
  D(i,j) = \min \begin{cases}
     D(i-1, j) + 1 & \text{deletion of } X[i]\\
     D(i, j-1) + 1 & \text{insertion of } Y[j]\\
     D(i-1, j-1) + \begin{cases} 2 & \text{if } X[i] \neq Y[j]\\ 0 & \text{if } X[i] = Y[j]\end{cases}
       & \text{substitution or copy}
  \end{cases}
\]
\textbf{Termination}: $D(n, m)$ is the distance.
\end{nbfact}

\begin{nbintuition}[reading the recurrence]
To turn $X[1..i]$ into $Y[1..j]$, look at the \emph{last} column of an optimal
alignment. Either the last character of $X$ was deleted (and the rest is an optimal
alignment of $X[1..i-1]$ with $Y[1..j]$), or the last character of $Y$ was inserted, or
the two last characters are aligned with each other (free if equal, a substitution
otherwise). The minimum over the three cases is the answer.
\end{nbintuition}

The general version (Jurafsky and Martin) uses arbitrary costs \texttt{del-cost},
\texttt{ins-cost}, \texttt{sub-cost}, which can depend on the letters, to model that some
letters are more likely to be inserted or confused than others; copying a letter costs
$\texttt{sub-cost}(x,x) = 0$.

\begin{nbwarning}[indices on the slide]
The slide defines $X$ of length $n$ and $Y$ of length $m$, but the recurrence loops over
$i = 1\dots M$, $j = 1\dots N$ and terminates with $D(N, M)$. The letters are swapped:
$i$ indexes $X$ and runs to $n$, $j$ indexes $Y$ and runs to $m$, and the result is
$D(n,m)$. The bottom-up loop also covers $1 \le i \le n$, $1 \le j \le m$ (the slide
writes strict inequalities).
\end{nbwarning}

\subsubsection{The table and the backtrace}
\slidesrc{Syntax: RE \& word: 81--82}

\begin{nbexample}{intention $\to$ execution}
Rows: the source \emph{intention}; columns: the target \emph{execution}; \texttt{\#} is
the empty prefix. Each cell is computed from its left, upper and upper-left neighbours.
\begin{center}
\small
\begin{tabular}{@{}c|cccccccccc@{}}
 & \# & \textbf{e} & \textbf{x} & \textbf{e} & \textbf{c} & \textbf{u} & \textbf{t} & \textbf{i} & \textbf{o} & \textbf{n} \\ \hline
\# & \cellcolor{nbteal!25}\textbf{0} & 1 & 2 & 3 & 4 & 5 & 6 & 7 & 8 & 9 \\
\textbf{i} & \cellcolor{nbteal!25}\textbf{1} & 2 & 3 & 4 & 5 & 6 & 7 & 6 & 7 & 8 \\
\textbf{n} & 2 & \cellcolor{nbteal!25}\textbf{3} & 4 & 5 & 6 & 7 & 8 & 7 & 8 & 7 \\
\textbf{t} & 3 & 4 & \cellcolor{nbteal!25}\textbf{5} & 6 & 7 & 8 & 7 & 8 & 9 & 8 \\
\textbf{e} & 4 & 3 & 4 & \cellcolor{nbteal!25}\textbf{5} & \cellcolor{nbteal!25}\textbf{6} & 7 & 8 & 9 & 10 & 9 \\
\textbf{n} & 5 & 4 & 5 & 6 & 7 & \cellcolor{nbteal!25}\textbf{8} & 9 & 10 & 11 & 10 \\
\textbf{t} & 6 & 5 & 6 & 7 & 8 & 9 & \cellcolor{nbteal!25}\textbf{8} & 9 & 10 & 11 \\
\textbf{i} & 7 & 6 & 7 & 8 & 9 & 10 & 9 & \cellcolor{nbteal!25}\textbf{8} & 9 & 10 \\
\textbf{o} & 8 & 7 & 8 & 9 & 10 & 11 & 10 & 9 & \cellcolor{nbteal!25}\textbf{8} & 9 \\
\textbf{n} & 9 & 8 & 9 & 10 & 11 & 12 & 11 & 10 & 9 & \cellcolor{nbteal!25}\textbf{8} \\
\end{tabular}
\captionof{table}{Levenshtein table. The bottom-right cell is the distance (8). Shaded
cells: one minimum-cost path found by the backtrace.}
\end{center}
Some cells, step by step:
\begin{align*}
  D(1,1) &= \min\bigl(D(0,1)+1,\ D(1,0)+1,\ D(0,0)+2\bigr) = \min(2, 2, 2) = 2\\
  D(4,1) &= \min\bigl(D(3,1)+1,\ D(4,0)+1,\ D(3,0)+0\bigr) = \min(5, 5, 3) = 3\\
  D(9,9) &= \min\bigl(D(8,9)+1,\ D(9,8)+1,\ D(8,8)+0\bigr) = \min(10, 10, 8) = 8
\end{align*}
In $D(1,1)$ the letters differ (\texttt{i}, \texttt{e}), so the diagonal costs 2; in
$D(4,1)$ and $D(9,9)$ they are equal, so the diagonal is a free copy.
\end{nbexample}

\textbf{Backtrace.} The table gives the distance, not the alignment. To get the
alignment, when filling each cell we remember \emph{which} neighbour(s) gave the minimum
(an arrow $\leftarrow$, $\uparrow$ or $\nwarrow$). After the table is full, we start from
the bottom-right cell and follow the arrows back to $(0,0)$: $\uparrow$ is a deletion,
$\leftarrow$ an insertion, $\nwarrow$ a substitution or a copy. The shaded path, read from
the top-left, is: delete \texttt{i}; \texttt{n}$\to$\texttt{e}; \texttt{t}$\to$\texttt{x};
copy \texttt{e}; insert \texttt{c}; \texttt{n}$\to$\texttt{u}; copy \texttt{t},
\texttt{i}, \texttt{o}, \texttt{n}. That is exactly the alignment shown above, with cost
$1 + 2 + 2 + 0 + 1 + 2 = 8$. When several neighbours tie, there are several optimal
alignments.

\begin{nbintuition}[cost of the algorithm]
The table has $(n+1)(m+1)$ cells, each computed in constant time from three neighbours,
so the algorithm takes $O(nm)$ time and space; the backtrace adds $O(n+m)$. The naive
search over all edit sequences is exponential.
\end{nbintuition}

\subsubsection{Exercises}
\notinslides

\begin{nbexercise}{1}{small tables}
Compute the Levenshtein table for \emph{cat} $\to$ \emph{cut} and for \emph{cat} $\to$
\emph{act}. Give the distance and an optimal alignment for each.
\end{nbexercise}
\begin{nbsolution}
\emph{cat} $\to$ \emph{cut} (rows c, a, t; columns c, u, t):
\[
\begin{array}{c|cccc}
   & \# & c & u & t\\ \hline
\# & 0 & 1 & 2 & 3\\
c  & 1 & 0 & 1 & 2\\
a  & 2 & 1 & 2 & 3\\
t  & 3 & 2 & 3 & 2
\end{array}
\]
Distance 2: copy \texttt{c}, substitute \texttt{a}$\to$\texttt{u} (cost 2), copy
\texttt{t}. Deleting \texttt{a} and inserting \texttt{u} also costs 2: there are two
optimal alignments.

\emph{cat} $\to$ \emph{act} (columns a, c, t):
\[
\begin{array}{c|cccc}
   & \# & a & c & t\\ \hline
\# & 0 & 1 & 2 & 3\\
c  & 1 & 2 & 1 & 2\\
a  & 2 & 1 & 2 & 3\\
t  & 3 & 2 & 3 & 2
\end{array}
\]
Distance 2, e.g.\ delete \texttt{c}, copy \texttt{a}, insert \texttt{c}, copy \texttt{t}.

\emph{Lesson}: swapping two adjacent letters costs 2. The Damerau--Levenshtein variant
adds \emph{transposition} as a fourth operation, because swaps are a very common typing
error.
\end{nbsolution}

\begin{nbsummary}[syntax at the word level]
\begin{itemize}[nosep]
  \item Word segmentation: spaces are only an approximation; Chinese, Japanese, Thai have
        no word spaces. MaxMatch = greedy longest dictionary word (good for Chinese,
        bad for English).
  \item Stop words: very frequent, little content; no universal list. A word frequent in
        a document may be topical: compare with the corpus.
  \item $\text{tf}(t,d)=\#t/\text{len}(d)$,
        $\text{idf}(t,C)=\log\bigl(\abs{C}/\abs{\{d: t\in d\}}\bigr)$,
        tf-idf = tf $\cdot$ idf; a word in every document has idf $=0$.
  \item Stemming = heuristic suffix chopping (Porter: measure $m$ in $[C](VC)^m[V]$,
        conditional rules in 7 steps; over- and under-stemming). Lemmatisation =
        dictionary + morphology, returns the lemma.
  \item Minimum edit distance: insertions, deletions, substitutions; Levenshtein:
        substitution costs 2. Dynamic programming: $D(i,0)=i$, $D(0,j)=j$,
        $D(i,j)$ = minimum over the upper neighbour $+1$ (deletion), the left
        neighbour $+1$ (insertion), the diagonal neighbour $+2$ (substitution) or $+0$
        (copy); the backtrace gives the alignment; time and space $O(nm)$.
\end{itemize}
\end{nbsummary}

